\section{Limitations and future work}
\label{sec:limitations}

\paragraph{Public benchmarks are short.}
The real test videos last seconds to minutes. The null footage is 12.6\,s for Ped2, 6.9\,min for Avenue and 26.9\,min for the Mivia-IWDD negative videos. Mivia-IWDD provides no per-camera calibration footage, so every test camera is unseen and calibration pools cameras. Pooling is one of three candidate explanations for the Mivia-IWDD PatchCore proportion of 9 of 50 normal videos (\cref{sec:real-data}), and per-camera calibration footage would test it directly. Long-horizon and per-camera validation therefore await operational multi-week data. The week-long simulation rests on diurnal and weather assumptions that such data would test.

\paragraph{Backbones are frozen and training-free.}
The backbone study uses frozen, training-free backbones because of the available compute budget; the experiments ran on a single cloud GPU. Detection power depends on the backbone. Optical flow detects 11 of 12 Ped2 episodes, while PatchCore detects 9 of 74 Avenue episodes and 10 of 50 Mivia-IWDD dumping videos, where flow detects 3 and none. Training state-of-the-art, per-dataset video anomaly detectors under the same leak-free, training-only protocol is future work and the natural next step for power. The Mivia-IWDD labels mark only the onset, so the detections cannot yet separate the dumping action from the waste left behind. A residue-specific model is also future work.

\paragraph{The theory assumes conditionally super-uniform p-values.}
\Cref{thm:episodic} requires the betting p-values to be conditionally super-uniform given the past. On real video the thinned p-values retain serial dependence. The lag-one correlation lies between 0.758 and 0.985 at the selected lags, and after a p-value below 0.1 the next p-value has mean 0.015 to 0.203, against 0.5 under the premise (\cref{tab:r3_dependence}). On real video the time-uniform bound is therefore used as an operating rule, and its realised per-video false-alarm rates are verified empirically. Extending the theory to dependent p-values, for example through block betting, is future work. In simulation, thinning by the decorrelation lag reduces the dependence but does not remove it, and the Ljung--Box test rejects independence for most cameras \citep{ljung1978}.

\paragraph{The shift study covers one family.}
The stricter support region, with slack 0.5, holds its level at every shift tested, and it was confirmed on fresh seeds. The evidence covers one activity-shift family, in one direction. For other shifts the abstention rate is the monitor we propose. The likelihood-ratio weighted path has no calibration-conditional proof, because its Kish effective-sample-size substitution is a heuristic that \cref{thm:beta} does not cover. The recalibration threshold of 0.05 on the abstention rate is a suggestion, not a validated rule.

\paragraph{The episodic advantage is configuration-specific.}
Episodes help when they are long, intermittent and recurring (\cref{tab:corner_sweep}). With 1600-frame episodes, $\pi=0.2$ and a 12{,}000-frame gap, the episodic mixture is ahead by 141.1\,s (95\% CI +79.6 to +210.2\,s). A single 800-frame episode favours the changepoint mixture by 55.7\,s (95\% CI $-111.2$ to $-14.0$\,s), dense and permanent changes are equivalent under the 1\,s margin of \cref{sec:episodes}, and both mixtures miss 200-frame episodes at $\alpha=0.01$. The delay comparison and the ablation use one configuration, four episodes in 60 minutes (\cref{sec:delay}). The guarantee does not depend on the episode model, but power does. Heavy-tailed episode durations, which dumping may produce, are untested.

\paragraph{Deployment considerations.}
The time-uniform guarantee is the horizon-free counterpart of the in-control average run length, and an alarm budget per monitoring period is its operational counterpart. Control charts also need a recalibration cadence, which statistical process control treats as a design choice \citep{montgomery2020introduction}, and an ROC operating point \citep{fawcett2006introduction} describes the trade-off that control rooms already use. Hourly alarm rates on real video are large, at 100.5 per hour of null footage for PatchCore on Mivia-IWDD, so an alarm budget should be set against them as well as against per-video rates. The choice of $\alpha$, against the costs of a false dispatch and of a missed event, is left to each operator. Every alarm should go to a named human reviewer before any action. A deployment also needs a retention policy and a route by which residents can contest a flagged event, because public-space surveillance raises accountability questions that the guarantee does not answer.
